\makeatletter
\let\thetitle\@title
\let\theauthor\@author
\makeatother

\newcommand{\togglepaper}[1][0]{
   \bibliography{../localbibliography}
   \papernote{\scriptsize\normalfont
     \theauthor.
     \titleTemp.
     To appear in:
     E. Di Tor \& Herr Rausgeberin (ed.).
     Booktitle in localcommands.tex.
     Berlin: Language Science Press. [preliminary page numbering]
   }
   \pagenumbering{roman}
   \setcounter{chapter}{#1}
   \addtocounter{chapter}{-1}
}

\newbool{bookcompile}
\booltrue{bookcompile}
\newcommand{\bookorchapter}[2]{\ifbool{bookcompile}{#1}{#2}}

% LLMの線
\DeclareTColorBox[auto counter, number within=chapter]
                 {LLMPrompt}{m O{0.8mm} O{\lsSeriesColor}}
  {
    graphical environment = tikzpicture,
    title = {#1},
    boxsep = 0pt,
    toptitle = 5mm,
    top = 5mm,
    bottom = 5mm,
    left = 5mm,
    right = 5mm,
    borderline = {#2}{0pt}{#3},
    beforeafter skip balanced = \baselineskip,
  }

\newcommand{\llmprompt}[1]{%
\begin{LLMPrompt}{LLM Prompt \thetcbcounter:}
#1
\end{LLMPrompt}
}

% \newcommand{\llmprompt}[1]{%
%   \vspace{\baselineskip} % 上に1行分の空白を追加
%   \noindent
%   \begin{tabular}{|p{.9\textwidth}|}
%     \hline
%     \footnotesize \textbf{LLM Prompt:}\\ % フォントサイズをさらに小さく
%     \footnotesize #1 \\
%     \hline
%   \end{tabular}
%   \\ % 改行を追加
%   \vspace{\baselineskip} % 下に1行分の空白を追加
% }

% LLM responseの線
\DeclareTColorBox[auto counter, number within=chapter]
                 {LLMResponse}{m O{black!12}}
      {
        graphical environment = tikzpicture,
        title = {#1},
        boxsep = 0pt,
        toptitle = 5mm,
        top = 5mm,
        bottom = 5mm,
        left = 5mm,
        right = 5mm,
        frame engine = path,
        extras = {frame engine=path},
        frame style = {fill=#2},
        sharp corners = all,
      }
  
\newcommand{\llmresponse}[1]{
\begin{LLMResponse}{LLM Response \thetcbcounter:}
#1
\end{LLMResponse}
}

% \newcommand{\llmresponse}[1]{
%   \vspace{\baselineskip} % 上に1行分の空白を追加
%   \noindent
%   \begin{tabular}{|p{\textwidth}|}
%     \hline
%     \footnotesize \textbf{LLM Response:}\\ % フォントサイズをさらに小さく
%     \footnotesize #1 \\
%     \hline
%   \end{tabular}
%   \\ % 改行を追加
%   \vspace{\baselineskip} % 下に1行分の空白を追加
% }

\renewcommand{\lsAcknowledgementTitle}{Acknowledgements}
